% WHAT A PAGE BREAK SPENDS.
%
% The penalty a page broke at is SPENT: \outputpenalty takes its value, and
% the node itself is left at ten thousand, so that it can never be a
% breakpoint again. The material after the break goes back to the
% contribution list with that node still in it, and an output routine that
% looks at the list sees the ten thousand:
%
%   \output={\showlists ...}
%   \vsize=1pt \hbox{}\penalty-10000
%
%   ### recent contributions:
%   \penalty 10000
%
% Without this the next page breaks at the same place for ever.
%
% HSTeX spent the penalty where the break was one it had WEIGHED and found
% best, and not where the break was FORCED -- a penalty of -10000 or worse
% fires the page builder from the node it is looking at, which has not been
% moved to the page yet, and that one was left as it stood.
\catcode`\{=1 \catcode`\}=2
\tracingonline=1 \showboxbreadth=8 \showboxdepth=8
\vsize=100pt \hsize=100pt \topskip=10pt
\message{[A a forced break]}
\output={\message{<forced \the\outputpenalty>}\showlists
  \global\setbox9=\box255}
\hbox{}\penalty-10000
\message{[B what the contribution list kept]}\showlists
\message{[C a break the builder weighed and chose]}
\output={\message{<weighed \the\outputpenalty>}\global\setbox9=\box255}
\vsize=20pt \hbox{}\penalty-5000 \hbox{}\hbox{}\hbox{}
\message{[D what the contribution list kept]}\showlists
\message{[done]}
\end
